\chapter{評価}
\label{chap:evaluation}

前の2つの章では、それぞれ独立した実験を提示し、各章で定量的・定性的な評価を行った。本章ではその両方から一歩引いて、プロジェクト全体を評価する。すなわち、結果が集合的に何を示しているか、証拠がどこで不十分か、そしてどのような限界が未解決のまま残っているかを論じる。

\section{実験全体にわたる結果のまとめ}

2つの実験（第~\ref{chap:development1_emo_steering}章および第~\ref{chap:development2_preverbal_vectors}章）は、大規模言語モデル（LLM）における情動的表現の構造と制御可能性について、相補的な問いに取り組んでいる。\par

第~\ref{chap:development1_emo_steering}章では、生のCAA方向ベクトルが共有された情動化成分 $\mathbf{g}$ に支配されていること、そしてそれを射影除去することで感情特有の残差 $\hat{\mathbf{r}}_e$ が得られることを示した。この残差はより識別力が高く、Plutchikの円環幾何学をより良く回復し、推論時に信頼性の高い情動的シフトをもたらす。共有成分は感情特有の信号を分離する上で分析的に有用だが、生成時の介入としては逆効果となる。\par

第~\ref{chap:development2_preverbal_vectors}章では、感情内残差部分空間は平坦ではないことを示した。名前付き感情はその内部次元数において大きく異なっており、プールドPCAにより2つの独立した抽出手法によって確認された3つの安定したメタ軸が回復される。これらの軸のうち2つ（$\mathbf{m}_1$、$\mathbf{m}_2$）は因果的に活性であり、感情ラベルを変えることなく情動的テクスチャのシフトをもたらす。このシフトはLLM評価者によって検証されており、一貫したステアリング区間を通じて意味の保持性は高く維持される。\par


\section{人間評価研究}
\label{sec:eval-human}

自動LLM評価者によって検出された情動的シフトが人間の読者にとって知覚可能であることを検証するため、付録~\ref{app:human_eval}に記載のフォームを使用した小規模な人間評価研究（$N=18$名）を実施した。本研究は、論文の3つの中心的主張に対応する3つのタスクで構成されている。

%TODO: Appendix 2: human eval form and instructions

\subsection{第1部：メタ軸テクスチャシフトの知覚可能性}

第1部では、$\mathbf{m}_1$によるステアリングで生じるテクスチャシフトが人間の読者に知覚可能かどうかを検証する。ここで焦点を当てるのは、感情カテゴリが同一のまま（例えば、joyはjoy、sadnessはsadnessのまま）であるときに、一方の文が他方と比べて人間によって
\begin{itemize}
    \item より身体的に圧倒的なものとして感じられる
    \item より直接的かつ本能的な方法で感情を表現している
    \item 緊迫感や生の強度を伝えている
\end{itemize}
かどうかである。\par

参加者は、同一の元発話から生成された一対の文を提示される。一方は\textit{興奮/反応的}な極方向にステアリングされた文（$\beta = +4$）、もう一方は\textit{穏やか/受容的}な極方向にステアリングされた文（$\beta = -4$）であり、どちらがどちらかは伝えられない。各ペアについて、参加者はどちらの文がより強烈で即時的な感情表現であるかを選択する。これは二者択一の強制選択（2AFC）タスクであり、参加者は中立の選択肢なしに、より強烈な文として必ずいずれか一方を選ばなければならない。仮説として、$\mathbf{m}_1$方向が感情内の\textit{興奮/即時的}なサブ軸を捉えているならば、人間評価者は$\beta=+4$の出力を比較対象よりも感情的に強烈かつ即時的であると一貫して選択するはずである。\par

\begin{table}[H]
\centering
\small
\begin{threeparttable}
\begin{tabular}{llcccc}
\toprule
\textbf{Q} & \textbf{感情} & \textbf{期待される} & \textbf{正答数} & \textbf{割合} & \textbf{$p$値} \\
 & & \textbf{方向} & \textbf{($N=18$)} & & \textbf{（二項検定）} \\
\midrule
Q1 & Joy     & 文~2 & 14 & 0.778 & $0.031^{*}$ \\
Q2 & Fear    & 文~1 & 14 & 0.778 & $0.031^{*}$ \\
Q3 & Disgust & 文~2 & 14 & 0.778 & $0.031^{*}$ \\
Q4 & Anger   & 文~1 & 17 & 0.944 & $<0.001^{***}$ \\
Q5 & Sadness & 文~2 & 17 & 0.944 & $<0.001^{***}$ \\
\midrule
\textit{全体} & --- & --- & 76 & 0.844 & $<0.001^{***}$ \\
\bottomrule
\end{tabular}
\begin{tablenotes}
\small
\item $^{*}p < 0.05$;\quad $^{**}p < 0.01$;\quad $^{***}p < 0.001$ （チャンスレベル $p_0 = 0.5$ に対する両側二項検定）。
\end{tablenotes}
\end{threeparttable}
\caption{第1部：ペアワイズ強度識別の結果（$N = 18$）。参加者は2つの文のうちどちらがより強い感情強度を示しているかを選択した。期待される方向は、$\beta=+4$（興奮極）にステアリングされた文を示す。}
\label{tab:human_eval_p1}
\end{table}

全90判断（$N=18$、各5ペア）において、期待される文は76回（$84.4\%$）選択され、$50\%$のチャンスレベルを大きく上回った。プールされた正確な二項検定によりこの効果が確認された（$p < 0.001$）。これにより、$\mathbf{m}_1$によるステアリングで生じるテクスチャシフトが意図した通りに人間の読者に知覚可能であることの強い証拠が得られた。特に\textit{anger}と\textit{sadness}では効果が顕著で、18名中17名が期待される文を選択した（各$p < 0.001$）。Joy、fear、disgustはやや低いものの依然として有意な一致を示した（$14/18$、$p = 0.031$）。\par

\subsection{第2部：感情カテゴリの保持}

第2部では、$\mathbf{m}_1$によるテクスチャシフトが知覚される感情カテゴリを変化させないことを検証する。第1部では参加者に強度の豊かな文を選ばせたが、第2部では各文に特定の感情ラベルを付与するよう求める。具体的には、$\beta$ を変化させても感情カテゴリが別のものに入れ替わるのではなく、同一カテゴリ内で表現スタイルがシフトすることを検証する。例えば、\textit{fear}において$\beta=-4$と$\beta=0$の両条件の文が\textit{fear}を表現していると判断されれば、$\beta$操作は感情そのものを変換するのではなく表現テクスチャを調整していることが示唆される。\par

\begin{table}[H]
\centering
\small
\begin{tabular}{llccc}
\toprule
\textbf{Q} & \textbf{目標感情} & \textbf{条件（$\beta$）} & \textbf{正答数（$N=18$）} & \textbf{割合（\%）} \\
\midrule
Q6  & Fear    & $\beta = -4$ & 17 & 94.4 \\
Q7  & Fear    & $\beta = 0$  & 17 & 94.4 \\
\midrule
Q8  & Anger   & $\beta = +4$ & 16 & 88.9 \\
Q9  & Anger   & $\beta = -4$ & 15 & 83.3 \\
\midrule
Q10 & Sadness & $\beta = +4$ & 17 & 94.4 \\
Q11 & Sadness & $\beta = -4$ & 17 & 94.4 \\
\midrule
\textit{全体} & --- & --- & 99 & 91.7 \\
\bottomrule
\end{tabular}
\caption{第2部：感情識別の結果（$N = 18$）。参加者はステアリングされた各文にPluchtikの8つの基本感情のいずれかをラベル付けした。$\beta$ は $\mathbf{m}_1$ 方向に適用された軸ステアリングの大きさを示す。}
\label{tab:human_eval_p2}
\end{table}

意図された感情とのカテゴリ一致率は全条件にわたって高い値を示した（表~\ref{tab:human_eval_p2}）。\textit{fear}では$\beta = -4$と$\beta = 0$の両条件で$94.4\%$、\textit{sadness}では$\beta = +4$と$\beta = -4$の両条件で$94.4\%$だった。\textit{anger}では$\beta = +4$で$88.9\%$（1名が\textit{disgust}、1名が\textit{どれも該当しない}と回答）、$\beta = -4$で$83.3\%$（2名が\textit{sadness}、1名が\textit{trust}と回答）だった。\par

感情ラベルは各感情における$\beta$値の変化に対して安定していた。全体として、108回答中99回答（$91.7\%$）が意図された感情と一致した。これらの結果は、$\mathbf{m}_1$ステアリングが感情カテゴリを置換することなく表現テクスチャをシフトさせることを確認しており、第~\ref{chap:development2_preverbal_vectors}章においてLLM評価者が報告した\texttt{subtle\_affective\_shift}スコアと一致している。\par

\subsection{第3部：プロトタイプステアリングの知覚可能性}

より強い目標感情ステアリングが目標感情強度の知覚可能な増加をもたらすかどうかを測定するため、第3部では参加者に一対の文を提示し、各文を独立して評価するよう求める。第1部と第2部が強制選択比較を用いるのに対し、第3部では各文に対して独立したリッカート評価を使用する。これにより2つの評価が切り離され、参加者はペアの中でより強烈なものを選ぶのではなく、各文の強度スコアをその内容に基づいて絶対的に付与する。\par

具体的には、各評価セットは元の発話、指定された目標感情、ステアリングなしの文（$\alpha_R = 0$）、および$\alpha_R = 64$でステアリングされた文から構成され、どの$\alpha_R$値で生成されたかは示されない。各文について、参加者は「この文において【目標感情】をどの程度強く感じますか？」を5点尺度（1 = \textit{全く感じない}；5 = \textit{非常に強く感じる}）で評価する。この尺度は知覚された目標感情強度を捉えるものであり、自然さや全体的な品質は問わない。\par

中心的仮説は、より強い目標感情ステアリングで生成された文は、ステアリングなしのベースラインよりも高い人間評価強度を得るはずである、すなわち6組の文ペアのそれぞれについて $\mu_{\alpha_R=64} > \mu_{\alpha_R=0}$ であるというものだ。\par

\begin{table}[H]
\centering
\small
\begin{threeparttable}
\begin{tabular}{llcccc}
\toprule
\textbf{セット} & \textbf{目標感情} & $\boldsymbol{\mu}_{\alpha=64}$ \textbf{（SD）} & $\boldsymbol{\mu}_{\alpha=0}$ \textbf{（SD）} & $\boldsymbol{\Delta}$ & $\boldsymbol{t(17)}$ \\
\midrule
Angerセット1 & Anger   & 3.56 (0.86) & 1.89 (0.90) & $+1.67$ & $6.52^{***}$ \\
Angerセット2 & Anger   & 3.61 (1.09) & 1.67 (0.69) & $+1.94$ & $8.26^{***}$ \\
Fearセット1  & Fear    & 3.72 (0.96) & 1.78 (0.81) & $+1.94$ & $6.81^{***}$ \\
Fearセット2  & Fear    & 3.89 (1.18) & 1.78 (0.73) & $+2.11$ & $8.30^{***}$ \\
Joyセット1   & Joy     & 4.17 (0.51) & 2.33 (0.97) & $+1.83$ & $7.08^{***}$ \\
Sadnessセット1 & Sadness & 4.39 (0.50) & 2.06 (0.42) & $+2.33$ & $14.43^{***}$ \\
\midrule
\textbf{全体} & & \textbf{3.89} & \textbf{1.92} & $+$\textbf{1.97} & --- \\
\bottomrule
\end{tabular}
\begin{tablenotes}
\small
\item $^{***}p < 0.001$ （両側対応$t$検定、$\mathit{df}=17$）。評価は1--5のリッカート尺度によるもの。
\end{tablenotes}
\end{threeparttable}
\caption{第3部：6つの感情・原文ペア（$N = 18$）にわたる、ステアリングあり（$\alpha_R = 64$）およびステアリングなし（$\alpha_R = 0$）の文に対する人間の平均強度評価。すべての差は期待される方向を示す。}
\label{tab:human_eval_p3}
\end{table}

ステアリングされた文はすべてのペアにおいてベースラインよりも高い平均評価を得ており（表~\ref{tab:human_eval_p3}）、全体的な平均差は $\Delta = +1.97$ ポイントだった。対応$t$検定により、個々のペアそれぞれについて効果が非常に有意であることが確認された（$t(17) \geq 6.52$、すべての場合において$p < 0.001$）。\par

LLM評価者の\texttt{target\_emotion\_intensity}スコアと人間の平均評価との間のSpearman相関係数は、12条件（6ペア $\times$ 2つの$\alpha_R$レベル）にわたって計算した結果 $\rho = 0.643$ であり、中程度の正の一致を示している。主な不一致は、$\alpha_R = 64$におけるFearセット~2で生じた。LLM評価者はスコア $0.00$ を付与した一方で、人間は同じ文を $3.89$ と評価しており、自動評価者が特定の表層的な構造においてfear強度を過小評価していることが示唆される。これらの結果はLLM評価者（LLM-as-judge）による評価を人間の知覚の代理指標として概ね支持しているが、対応は不完全であり、人間によるキャリブレーションが依然として重要である。\par

以上の3つのタスクを総合すると、提案されたステアリング方向が人間の読者によって確実に知覚可能な情動的差異に対応していることが示されており、自動評価の結果を裏付け、第~\ref{sec:eval-limitations}節で指摘したLLM評価者の限界を部分的に解消している。
\section{限界}
\label{sec:eval-limitations}

\subsection{LLM評価者による評価}

両実験では\texttt{GPT-4o}または\texttt{GPT-4o-mini}を主要な評価者として使用している~\cite{gu2025surveyllmasajudge}。この設定の一つの限界は、評価者とステアリング対象モデル（\texttt{Llama-3.1-8B-Instruct}）の両方が重複するコーパスで指示チューニングされているため、評価者が独立した情動的信号を評価するのではなく、評価対象となっているのと同じ表層レベルの文体的規則性を符号化している可能性があることだ。\par

第2の懸念点はラベルの網羅性である。評価者はPluchtikの8つの感情ラベルとプロンプトで提供された軸極の説明に基づいて出力を評価するため、このラベル空間に収まらない情動的シフトは誤ってスコア付けされるか、最も近いラベルに集約されてしまう可能性がある。\par

第3の懸念点は外的妥当性である。評価者スコアは人間の知覚の連続的な代理指標であり、人間による検証研究（$N = 18$）は基準となるキャリブレーションを確立するには規模が小さすぎる。したがってこれらの知見は、LLMが知覚した情動的シフトとして解釈されるべきである。それらが人間の読者に知覚されるシフトに対応するかどうかは、より大規模な研究で確立される必要がある。\par

\subsection{単一層介入}

すべてのステアリング実験では、第~\ref{sec:layer_selection}節の層診断に従い、単一のトランスフォーマー層（第13層）にベクトルを注入している。これは原理に基づいた選択だが、単純化でもある。情動的処理が単一の層に局在している必要はなく、異なる次元は異なる深さで最もよくステアリングされる可能性がある。$\mathbf{m}_3$（対人的関与）に観察されたほぼゼロのステアリング効果はこの可能性と一致している。その軸は単一サイトの注入では効果的に摂動できない方法で層をまたいで分布している可能性がある。本研究では、多層介入および軸ごとの層選択は検討されなかった。注入深度を軸ごとに変化させる体系的な研究が、$\mathbf{m}_3$のヌル効果が表現の幾何学的構造を反映しているのか、それとも単に介入サイトとその軸がエンコードされている場所とのミスマッチによるものなのかを判断するために必要である。\par

\subsection{感情の範囲と一般化可能性}

CAAベクトルはPluchtikの8つの基本感情から引き出された対比ペアによって定義されており、この選択が結果として得られる表現空間の幾何学的構造を制約している。この分類体系は理論的なコミットメントである。他の枠組み（例えば、Russellの価値-覚醒度円環モデル、または音楽誘発状態のためのGEMS）は感情空間を異なる方法で分割し、異なる対比ペアを生み出し、結果として異なるCAA方向をもたらすだろう。分解 $\mathbf{g}$、$\hat{\mathbf{r}}_e$、$\mathbf{m}_k$ がPluchtikのカテゴリに特有なものなのか、それとも情動的幾何学のより一般的な性質を反映しているのかは未解決の問いである。\par

さらなる交絡要因がデータ生成手順によってもたらされる。書き換えデータセットはGPT-4.1-miniを使用して英語対話コーパスから構築されている。残差ストリームで捕捉された情動的シフトは、\texttt{Llama-3.1-8B}の固有の情動的幾何学ではなく、GPT-4.1-miniの文体的規則性を部分的に反映している可能性がある。生成モデルの寄与を目標モデルの内部表現から切り離すには、書き換えモデルを変化させながら目標モデルを固定する対照研究が必要である。\par

すべての実験は単一のモデルファミリー（\texttt{Llama-3.1-8B-Instruct}）を使用している。この分解が他のアーキテクチャやモデルスケールに一般化するかどうかは未検証である。\par

\subsection{メタ軸ラベルの解釈}

$\mathbf{m}_1$--$\mathbf{m}_5$ に付与されたラベル（表~\ref{tab:meta_axis_interpretations}）は、LLM評価者に極一致した例の対比を特徴付けるよう促すことで生成された事後的解釈だ。このようなラベルは、潜在的な情動的構造よりも生成されたテキストの表層的規則性（トピック、レジスター、表現）を部分的に追跡している可能性があり、確定的なアノテーションではなく妥当な特徴付けとして扱われるべきである。\par

\section{批判的評価}
\label{sec:eval-appraisal}

実験は情動的ステアリング方向の提案された分解に対して相当な証拠を提供しているが、その証拠の強さは主張によって異なる。本節では、実験によって直接支持される知見と、より慎重に扱われるべき解釈を区別する。\par

最も強力な実証的主張は、感情関連のCAA方向が大きな共有成分と小さな感情特有の残差を含むというものだ。いくつかの収束した結果がこれを支持している。生のCAAベクトルは互いに非常に類似しており、信号の多くが特定の名前付き感情をエンコードするのではなく、中立的な言語から感情的な言語への一般的なシフトを反映していることを示唆している。共有方向 $\mathbf{g}$ を除去した後、残差ベクトル $\hat{\mathbf{r}}_e$ は感情カテゴリを区別するためにより有用になり、Plutchik感情間の期待される関係をより良く反映する。これらの残差方向は因果的にも活性である。生成時に注入すると、予測可能な目標感情シフトをもたらす。したがって $\hat{\mathbf{r}}_e$ は単なる記述的なアーティファクトではなく、機能的なステアリング方向である。\par

第2の十分に支持された主張は、情動的ステアリングが離散的な感情ラベルの選択によって尽くされるわけではないというものである。第~\ref{chap:development2_preverbal_vectors}章は、各名前付き感情の内部に意味のある構造が存在することを示している。メタ軸 $\mathbf{m}_1$ と $\mathbf{m}_2$ は、感情カテゴリを大きく保持したまま感情の表現テクスチャを調整する。人間評価はこの点を強化しており、参加者は $\mathbf{m}_1$ シフトを強度と即時性の差として知覚しながら、ほとんどの場合で同一の意図された感情ラベルを付与した。これはモデルが\emph{どの}感情が表現されているかだけでなく\emph{どのように}表現されているかも符号化していることを示唆している。\emph{前言語的}という用語は狭い技術的意味で使用されている。$\mathbf{m}_k$ はサブラベルであり、前言語的ではない。モデルには現象的な経験はなく、回復されるのは心理学的に根拠のある情動的次元と相関する活性化幾何学である。\par

LLM評価者による評価の使用には注意が必要である。自動評価者は多数の生成物と複数のステアリング強度にわたる一貫した評価を可能にし、人間評価は特に $\mathbf{m}_1$ の知覚可能性とプロトタイプステアリング下での目標感情強度の増加について自動評価結果を概ね支持している。しかし、LLM評価者スコアと人間評価との相関は中程度にとどまり、fear条件では明確な乖離が生じている。したがってLLM評価者は有用なスケーラブルな代理指標であるが、人間評価の完全な代替ではない。本論文の最も強力な主張は、それに応じて両方の証拠源によって支持されたものである。\par

実験設計は意図的に制御されている。貪欲法によるデコード、ホールドアウトされた原文の分割、固定された層、および構造化された評価プロンプトにより、ノイズが低減されてベクトルの因果効果が分離される。この制御は強みであると同時に限界でもある。実験は対話的な原文発話からの単文書き換えを評価しており、同一の分解が長文パッセージ、マルチターン会話、多言語設定、または他のジャンルに転用可能かどうかは未確立である。したがって結果は、制約された設定内での信頼性の高いメカニズムの証拠として理解されるべきであり、普遍的な情動制御の証明としてではない。\par

総じて、実験は限定的な主張を支持している。研究対象モデルにおける情動的生成は、共有された情動化方向、感情特有の残差方向、および感情カテゴリ内で表現的な性格を調整するサブラベルテクスチャ軸に分解できる。これらの成分は幾何学的に意味があり、因果的に活性であり、人間の読者に知覚可能である。同一の分解がモデルファミリー、スケール、または学習レジームにわたって安定しているかどうかは、単一モデル・単一層の設計では解決できない未解決の問いである。\par
